% =============================================================================
% Capítulo 1 — Introdução                                         (peso: leve)
% -----------------------------------------------------------------------------
% Narrativa-base: pitch do plano (§11); lacuna (§2.1); problema (§3);
% objetivos (§4, reescritos após o feedback do orientador "parece
% metodologia": métricas e protocolo saíram dos específicos e ficam no Cap. 4).
% Ressalva A28: enquanto aberto, NÃO afirmar que as responsabilidades do harness
% já foram mapeadas contra R21.
% Números de contexto (Stack Overflow, Menlo, IEA): relatórios, acesso 09/10/2026.
% =============================================================================
\chapter{Introdução}
\label{cap:introducao}

\section{Contexto e motivação}
\label{sec:contexto}

Os modelos de linguagem de grande escala (\textit{Large Language Models},
LLMs)\glsadd{LLM} deixaram de ser objeto restrito de pesquisa e passaram a
compor o cotidiano do desenvolvimento de \textit{software}. Na pesquisa anual
da plataforma Stack Overflow, a proporção de desenvolvedores que usam ou
pretendem usar ferramentas de inteligência artificial (IA)\glsadd{IA} subiu de
76\% em 2024 para 84\% em 2025, e 31\% dos respondentes já declaravam usar
agentes de IA \cite{stackoverflow2025}. O movimento não se limita às
ferramentas de apoio ao programador: os próprios produtos passaram a incorporar
modelos. Segundo levantamento da Menlo Ventures com líderes técnicos de empresas
e \textit{startups}, o gasto corporativo com interfaces de programação de
aplicações (\textit{Application Programming Interfaces},
APIs)\glsadd{API} de LLMs mais que dobrou em seis meses, de 3,5 bilhões de
dólares em novembro de 2024 para 8,4 bilhões em meados de 2025
\cite{menlo2025}.

Em uma aplicação \textit{web}, essa incorporação costuma ocorrer da forma mais
direta possível: o código de negócio importa o kit de desenvolvimento
(\textit{Software Development Kit}, SDK)\glsadd{SDK} do provedor escolhido e
chama o modelo nos pontos em que precisa dele. Cada provedor, porém, define o
próprio formato de mensagens, de declaração e de chamada de ferramentas, de
erros e de limites de uso. À medida que as integrações deixam de ser uma
pergunta isolada ao modelo e passam a envolver a execução de ferramentas
externas em ciclos de várias etapas --- o que \citeonline{sapkota2026} chamam de
agentes de IA, em contraste com sistemas multiagentes de \textit{IA agêntica}
---, a quantidade de código dependente dessas convenções cresce. Esse tipo de
dependência já foi descrito como dívida técnica em sistemas de aprendizado de
máquina: o \textit{glue code} que adapta a aplicação à API de um pacote
específico tende a prender o sistema às particularidades desse pacote e a
encarecer qualquer substituição posterior \cite{sculley2015}.

A Engenharia de Software conhece há décadas a resposta de princípio para esse
problema. \citeonline{parnas1972} propôs que as decisões de projeto com maior
probabilidade de mudar fossem encapsuladas em módulos que as escondessem do
restante do sistema, e a modificabilidade é hoje tratada como um atributo de
qualidade arquitetural a ser deliberadamente projetado \cite{bass2021}. Fora do
domínio da IA, a dependência de um único fornecedor (\textit{vendor lock-in})
é um problema arquitetural consolidado: a literatura de computação em nuvem
reúne um amplo conjunto de abordagens de interoperabilidade e portabilidade
baseadas em camadas intermediárias \cite{kaur2017}, a noção de
\textit{middleware} como camada reutilizável que absorve a heterogeneidade entre
plataformas está bem estabelecida \cite{issarny2007}, e há comparações empíricas
entre funções \textit{serverless} escritas sobre bibliotecas de abstração e
sobre as APIs nativas do provedor \cite{mo2023}.

No caso dos modelos de IA, a escolha do provedor é justamente uma decisão com
alta probabilidade de mudar. O mesmo levantamento da Menlo Ventures registra
que a participação da OpenAI no uso corporativo de LLMs caiu de 50\%, ao fim de
2023, para 25\% em meados de 2025, enquanto a Anthropic chegou a 32\% e o
Google a 20\% \cite{menlo2025}. Ainda assim, apenas 11\% das equipes
consultadas trocaram de fornecedor no último ano, ao passo que 66\% apenas
atualizaram o modelo dentro do provedor que já utilizavam \cite{menlo2025}.
Esses números não isolam a causa da baixa taxa de troca --- desempenho,
contratos e preço também pesam ---, mas são coerentes com a hipótese de que
trocar de provedor tem custo técnico relevante para quem integrou o modelo
diretamente ao código.

Uma alternativa arquitetural é interpor, entre a aplicação e o provedor, uma
camada dedicada à integração com os modelos. Na prática recente, o termo
\harness{} passou a designar a camada de \textit{software} que envolve o modelo
e se responsabiliza pela execução de ferramentas, pelo gerenciamento de estado
e pela recuperação de erros. Neste trabalho, \harness{} é definido
operacionalmente como a \textbf{camada arquitetural intermediária entre a
aplicação \textit{web} e o provedor de IA}, que abstrai e centraliza as
interações da aplicação com modelos e agentes, de modo que a aplicação não
acessa diretamente os SDKs dos provedores. Atribuem-se a essa camada a
padronização da comunicação com os provedores, a orquestração de chamadas e a
execução controlada de ferramentas, o gerenciamento de contexto e estado, o
tratamento de falhas e o registro estruturado de eventos; ela não implementa
regras de negócio nem treina modelos. Essas responsabilidades dialogam com o
catálogo de padrões arquiteturais para agentes baseados em modelos de fundação
proposto por \citeonline{liu2025}, que serve de referência para o projeto da
camada.

Essa camada, contudo, não é gratuita. Cada nível de indireção acrescenta
trabalho a toda requisição --- normalização de mensagens, serialização,
orquestração do ciclo de ferramentas, registro de eventos ---, e esse trabalho
consome energia. A questão ganhou peso com a escala atual da IA. A Agência
Internacional de Energia estima que os centros de dados consumiram cerca de
415~TWh em 2024, aproximadamente 1,5\% da eletricidade mundial, e projeta que
esse consumo mais que dobre até 2030, chegando a cerca de 945~TWh, com a IA como
principal vetor de crescimento \cite{iea2025}. A área de \textit{Green AI}
propõe tratar a eficiência energética como critério de avaliação ao lado da
acurácia \cite{schwartz2020}, depois que \citeonline{strubell2019} chamaram
atenção para o custo energético e ambiental do treinamento de modelos de
linguagem.

Na fase de uso, que se repete a cada requisição, os números também são
expressivos. \citeonline{luccioni2024} mostram que arquiteturas generativas de
propósito geral chegam a ser ordens de grandeza mais custosas que sistemas
específicos para a mesma tarefa: gerar texto consumiu em média 0,047~kWh a cada
mil inferências, contra 0,002~kWh da classificação de texto. Além disso, o
consumo depende fortemente da pilha de \textit{software}, e não apenas do
modelo: \citeonline{fernandez2025} verificaram que estimativas baseadas em
contagem de operações de ponto flutuante subestimam o consumo real e que
otimizações bem aplicadas reduzem a energia total em até 73\%.
\citeonline{vellaisamy2026} decompõem a sobrecarga de orquestração executada no
processador hospedeiro, fora da unidade de processamento gráfico
(GPU)\glsadd{GPU}, e mostram que ela é um custo oculto da inferência, enquanto
\citeonline{ifath2026} caracterizam os compromissos entre desempenho e energia
em fluxos agênticos de múltiplas requisições, o mesmo padrão de uso que uma
camada de \harness{} orquestra. Decorre daí que uma decisão puramente
arquitetural, como acrescentar uma camada entre a aplicação e o modelo, pode
ter consequências energéticas mensuráveis, que raramente entram na avaliação
dessa decisão.

\section{Problema e lacuna de pesquisa}
\label{sec:problema}

A literatura já propõe arquiteturas para integrar modelos e agentes a sistemas
de \textit{software} e discute os compromissos envolvidos. O catálogo de
\citeonline{liu2025} descreve, para cada padrão, contexto, forças e
\textit{trade-offs}, mas o faz qualitativamente, a partir de revisão de
literatura. \citeonline{kandogan2025} propõem uma arquitetura de referência
para IA composta, com registros de agentes e de dados que abstraem provedores e
APIs proprietárias, ilustrada por um caso de uso, sem comparação controlada com
uma integração direta. \citeonline{zhu2026}, ao comparar
\textit{frameworks} de orquestração de agentes baseados em LLMs, destacam que
parte relevante das alegações sobre esses \textit{frameworks} provém dos
próprios fornecedores, e não de avaliação revisada por pares.

A lacuna, portanto, não está na ausência de discussão sobre os
\textit{trade-offs} dessa camada, mas na ausência de \textbf{medição
controlada} deles: não foram encontrados estudos que comparem, sob as mesmas
condições, uma aplicação com e sem a camada de integração, quantificando ao
mesmo tempo o benefício mais alegado da camada --- a facilidade de trocar de
provedor --- e o preço que ela cobra em energia. Sem essa medição, a decisão de
adotar ou não uma camada de \harness{} fica apoiada em argumentos de princípio
ou em alegações de fornecedores.

Diante disso, este trabalho busca responder à seguinte questão de pesquisa:
\textit{a introdução de uma camada arquitetural intermediária entre uma
aplicação \textit{web} e um provedor de IA reduz o custo de substituição do
provedor, e qual é o custo energético dessa redução?}

Parte-se da hipótese de que a camada reduz o custo de substituição e cobra por
isso um custo energético; essa hipótese é tratada como algo a testar, e não
como pressuposto. Caso a camada reduza o custo de substituição sem acréscimo
relevante de energia, ou caso acrescente custo energético sem reduzir o esforço
de substituição, ambos os resultados são informativos para quem precisa tomar a
decisão arquitetural.

\section{Justificativa}
\label{sec:justificativa}

Do ponto de vista prático, desenvolvedores e arquitetos que integram LLMs a
aplicações \textit{web} decidem, explícita ou implicitamente, entre acoplar o
código ao SDK de um provedor ou investir em uma camada de abstração, própria ou
oferecida por um \textit{framework}. Em um mercado no qual a liderança entre
provedores mudou em menos de dois anos \cite{menlo2025}, o custo de trocar de
provedor deixa de ser hipotético; e, em um cenário de demanda crescente por
energia nos centros de dados \cite{iea2025}, qualquer sobrecarga paga a cada
requisição é multiplicada pelo volume de uso. Evidência quantitativa sobre os
dois lados dessa decisão permite tomá-la com base em dados, e não apenas em
princípios de projeto ou em material promocional.

Do ponto de vista científico, o trabalho contribui para a arquitetura de
\textit{software} de sistemas com IA ao submeter a avaliação experimental uma
decisão que, até aqui, foi discutida principalmente de forma qualitativa
\cite{liu2025}, e aproxima a Engenharia de Software da agenda de \textit{Green
AI} \cite{schwartz2020} ao tratar o consumo de energia como atributo de
qualidade de uma decisão arquitetural, e não apenas do modelo. O estudo segue
os princípios de experimentação controlada em Engenharia de Software
sistematizados por \citeonline{wohlin2024}, comparando duas versões
funcionalmente equivalentes de uma mesma aplicação: uma integrada ao provedor
por meio da camada e outra integrada diretamente.

O recorte em duas dimensões é deliberado. A substituibilidade do provedor
representa o benefício mais associado à camada, e o custo energético, o preço
que ela cobra; juntos, formam um compromisso fechado sobre a mesma decisão.
Outros atributos usualmente atribuídos a uma camada desse tipo, como a
recuperação de falhas e a auditabilidade, serão implementados por serem
pré-requisitos técnicos do experimento, mas não serão avaliados. Por isso, a
conclusão do trabalho se limita a responder se a camada compensa pelo
desacoplamento que proporciona, e não se ela compensa de modo geral.

Por fim, o trabalho é viável no prazo de um trabalho de conclusão de curso. A
medição energética apoia-se em protocolo de telemetria de \textit{hardware} já
empregado pelo autor em trabalho anterior \cite{barros2026}, e o experimento
utiliza infraestrutura com GPU disponível na instituição e modelos abertos
executados localmente.

\section{Objetivos}
\label{sec:objetivos}

\subsection{Objetivo geral}
\label{subsec:objetivo-geral}

Avaliar experimentalmente o \textit{trade-off} entre desacoplamento e eficiência
energética introduzido por uma camada arquitetural intermediária
(\harness{}) entre aplicações \textit{web} e provedores de modelos de IA, em
comparação com a integração direta ao provedor.

\subsection{Objetivos específicos}
\label{subsec:objetivos-especificos}

\begin{alineas}
  \item projetar uma camada de \harness{} para a integração de modelos de IA em
        aplicações \textit{web}, delimitando suas responsabilidades com base em
        padrões arquiteturais publicados;
  \item quantificar o custo de substituição do provedor de IA em uma aplicação
        \textit{web} com e sem a camada de \harness{};
  \item mensurar o custo energético que a camada de \harness{} acrescenta à
        execução de requisições simples e agênticas a um modelo de IA;
  \item identificar em que condições a redução do custo de substituição
        compensa o custo energético da camada.
\end{alineas}

% \section{Estrutura do trabalho}
% \label{sec:estrutura}
% Fonte: plano §8. Redigir quando os demais capítulos estiverem estáveis
% (um parágrafo por capítulo, com \ref{cap:...}).
